% BHU PYQ -- Professional Exam Template
\documentclass[11pt,a4paper]{article}

\usepackage[a4paper,top=2.5cm,bottom=2.5cm,left=2.8cm,right=2.8cm,headheight=14pt]{geometry}
\usepackage{amsmath,amssymb}
\usepackage[T1]{fontenc}
\usepackage[utf8]{inputenc}
\usepackage{lmodern}
\usepackage{microtype}
\usepackage{fancyhdr}
\usepackage{enumitem}
\usepackage{hyperref}
\usepackage{lastpage}
\usepackage{parskip}

\hypersetup{colorlinks=true,urlcolor=black,linkcolor=black,
  pdftitle={B.Sc. (Hons.) Zoology Sem-V ZOB-502 -- BHU PYQ}}

\pagestyle{fancy}
\fancyhf{}
\renewcommand{\headrulewidth}{0.4pt}
\renewcommand{\footrulewidth}{0.4pt}
\fancyhead[L]{\small   \textit{Banaras Hindu University}}
\fancyhead[R]{\small   \textit{Zoology Paper: ZOB-502}}
\fancyfoot[C]{\small Page  \thepage\ of \pageref{LastPage}}


\newlist{parts}{enumerate}{3}
\setlist[parts,1]{label=(\alph*),leftmargin=*,itemsep=4pt,topsep=3pt}
\setlist[parts,2]{label=(\roman*),leftmargin=*,itemsep=2pt,topsep=2pt}
\setlist[parts,3]{label=(\arabic*),leftmargin=*,itemsep=2pt,topsep=2pt}

\newcommand{\pts}[1]{\hfill   \textit{[#1]}}


\begin{document}


\begin{center}
  {\large\bfseries BANARAS HINDU UNIVERSITY}\\[2pt]
  {\normalsize B.Sc. (Hons.) Semester V Examination, 2016--17}\\[6pt]
  \rule{0.8\linewidth}{0.5pt}\\[6pt]
  {\large\bfseries Zoology}\\[4pt]
  {\large\bfseries Paper No. ZOB-502 : Functional Anatomy of Chordates}\\[4pt]
  {\normalsize Time: 3 hours \hfill Full Marks: 70}
\end{center}

\rule{\linewidth}{0.4pt}

\smallskip



\noindent   \textit{Note: Write your Roll No. at the top immediately on the receipt of this question paper. The figures in the right-hand margin indicate marks. Answer six questions including Question 1 which is compulsory.}

\bigskip



\noindent   \textbf{Question 1.} Answer the following parts:

\begin{parts}
  \item State whether the following statements are True or False. Justify your answers in 1 or 2 sentences each:\pts{$4  \times 2 = 8$}
  
\begin{parts}
    \item Tree like arborescent organs act as accessory respiratory organ in   \textit{Anabas}.
    \item Ceruminous glands are modified Sudoriferous glands.
    \item In Agnathans gill system also function as a food-collecting system.
    \item Dermatocranium forms the jaws in gnathostomes.
  \end{parts}

  \item Define the following:\pts{$3  \times 1 = 3$}
  
\begin{parts}
    \item Amphicoelous vertebra
    \item Mullerian duct
    \item Parabronchi.
  \end{parts}

  \item Match the following:\pts{$6  \times \frac{1}{2} = 3$}
  
\begin{center}
  
\begin{tabular}{llp{1.5cm}ll}
    \hline
    \multicolumn{2}{c}{   \textbf{Column--A}} & & \multicolumn{2}{c}{   \textbf{Column--B}} \\
    \hline
    (i) & Meckel's cartilage & & (1) & Solenodont condition \\
    (ii) & Contour feather & & (2) & Poison gland \\
    (iii) & Ruminants & & (3) & Ganoid scales \\
    (iv) & Parotid gland & & (4) & Pyloric caeca \\
    (v) & Primitive bony fish & & (5) & Quills \\
    (vi) & Ray finned fish & & (6) & Lower jaw \\
    & & & (7) & Plumule \\
    & & & (8) & Sebaceous gland \\
    \hline
  \end{tabular}
  \end{center}
  \medskip

  \item Differentiate between the following:\pts{$3  \times 2 = 6$}
  
\begin{parts}
    \item Horns and Antlers
    \item Secodont condition and Lophodont condition
    \item Bipartite and Bicornuate uterus.
  \end{parts}
\end{parts}

\pagebreak

\medskip


\noindent   \textbf{Question 2.} With the help of neat labelled diagrams, explain different types of jaw suspensorium in vertebrates.\pts{10}

\medskip


\noindent   \textbf{Question 3.} Give an illustrated account of gill system in shark.\pts{10}

\medskip


\noindent   \textbf{Question 4.} What are aortic arches? Give an account of modifications in aortic arches from elasmobranchs to anurans with the help of neat labelled diagrams.\pts{10}

\medskip


\noindent   \textbf{Question 5.} Explain the evolution of cerebellum in vertebrates.\pts{10}

\medskip


\noindent   \textbf{Question 6.} Explain the following in brief:\pts{$2  \times 5 = 10$}

\begin{parts}
  \item Dentition in mammals
  \item Hepatic portal system.
\end{parts}

\medskip


\noindent   \textbf{Question 7.} Give an account of different types of kidney tubules in vertebrates.\pts{10}

\medskip


\noindent   \textbf{Question 8.} Write short notes on the following:\pts{$2  \times 5 = 10$}

\begin{parts}
  \item Digital cornifications
  \item Mesonephros kidney.
\end{parts}

\vfill
\rule{\linewidth}{0.4pt}

\begin{center}\small
\href{https://bhu-exam.vercel.app/}{Click Here}\\[4pt]\href{https://me-aryan.vercel.app/}{Click Here}\\[4pt]\href{https://quashan-me.vercel.app/}{Click Here}
\end{center}

\end{document}
